\documentclass[x11names]{beamer}
\usepackage{etex,beamerfoils}
\usepackage{array,pifont,amsmath,latexsym,epsfig,amsfonts,yfonts}
\usepackage{amstext,amsopn,amsxtra,amsgen,amsbsy,amscd,amssymb,dcolumn,graphicx,setspace,pgfplots,booktabs,natbib}
\usepackage[normalem]{ulem}
\usepackage[all]{xy}
\beamertemplatetheoremsnumbered
\setbeamertemplate{caption}[numbered]
\usetheme{Copenhagen}
\definecolor{NewBlue}{RGB}{4, 114, 155}
\definecolor{magenta}{cmyk}{0,1,0,0}
\usecolortheme[named=NewBlue]{structure}
%\usecolortheme[named=SkyBlue4]{structure}

\newcommand{\D}{\displaystyle}
\newcommand{\E}{\begin{eqnarray*}}
\newcommand{\F}{\end{eqnarray*}}
\newcommand{\EE}{\begin{eqnarray}}
\newcommand{\FF}{\end{eqnarray}}
\newcommand{\IZ}{\begin{itemize}}
\newcommand{\ZI}{\end{itemize}}
\newcommand{\EN}{\begin{enumerate}}
\newcommand{\NE}{\end{enumerate}}
\newcommand{\itemc}{\item[$\circ$]}
\newcommand{\IZdash}{\begin{itemize} \renewcommand{\labelitemi}{-}}
\newcommand{\tn}{\textnormal}

% New commands to assist Beamer
\newcommand{\itemi}{\item<1->}
\newcommand{\itemii}{\item<2->}
\newcommand{\itemiii}{\item<3->}
\newcommand{\itemiv}{\item<4->}
\newcommand{\itemv}{\item<5->}
\newcommand{\itemvi}{\item<6->}
\newcommand{\itemvii}{\item<7->}
\newcommand{\itemviii}{\item<8->}
\newcommand{\itemix}{\item<9->}
\newcommand{\itemx}{\item<10->}

\newcommand{\itemci}{\item[$\circ$]<1->}
\newcommand{\itemcii}{\item[$\circ$]<2->}
\newcommand{\itemciii}{\item[$\circ$]<3->}
\newcommand{\itemciv}{\item[$\circ$]<4->}
\newcommand{\itemcv}{\item[$\circ$]<5->}
\newcommand{\itemcvi}{\item[$\circ$]<6->}
\newcommand{\itemcvii}{\item[$\circ$]<7->}
\newcommand{\itemcviii}{\item[$\circ$]<8->}
\newcommand{\itemcix}{\item[$\circ$]<9->}
\newcommand{\itemcx}{\item[$\circ$]<10->}


\newcommand{\fr}{\frame}
\newcommand{\ft}{\frametitle}
\newcommand{\ons}{\onslide}
\newcommand{\tc}{\textcolor}

\pgfplotsset{height =7cm,compat=1.7,width=7cm,
    standard/.style={
    	unit vector ratio*=1 1 1,
        axis x line=middle,
        axis y line=middle,
        enlarge x limits=0.15,
        enlarge y limits=0.15,
        every axis x label/.style={at={(current axis.right of origin)},anchor=north west},
        every axis y label/.style={at={(current axis.above origin)},anchor=north east}
    }
}
\newcommand{\drawge}{-- (rel axis cs:1,0) -- (rel axis cs:1,1) -- (rel axis cs:0,1) \closedcycle}
\newcommand{\drawle}{-- (rel axis cs:1,1) -- (rel axis cs:1,0) -- (rel axis cs:0,0) \closedcycle}

\usetikzlibrary{patterns}
\usetikzlibrary{arrows}
\usetikzlibrary{shapes}

\renewcommand{\bibsection}{}

% Watermark example: https://texblog.org/2012/02/17/watermarks-draft-review-approved-confidential/
\usepackage{draftwatermark}
\SetWatermarkText{This content is protected and may not be shared uploaded or distributed.}
\SetWatermarkScale{0.25}
\SetWatermarkAngle{0}
\SetWatermarkColor[gray]{1.0}
\setbeamercolor{background canvas}{bg=}%transparent canvas

\title[Dynamic Optimization Intro\hspace{2em}\insertframenumber/
\inserttotalframenumber]{Introduction to Dynamic Optimization}
\author{Brian C.~Jenkins \\ ~ \\University of California, Irvine}

\begin{document}

\frame{\maketitle}



\fr{
\ft{Dynamic Optimization}

\IZ
\itemi Dynamic optimization is at the heart of macroeconomic theory. 

\

\itemii People and firms routinely make decisions that affect their future opportunities. 

\

\itemiii For example, a household chooses how much of its income to consume today and to save for the future.
\ZI
}


\fr{
\ft{A Two-Period Cake-Eating Problem}

\IZ
\item A person has some initial quantity of cake

\

\itemii The person receives utility from consuming cake in period 0 and in period 1 only.

\

\itemiii Problem: what is the optimal cake consumption path through periods 0 and 1?
\ZI
}


\fr{
\ft{Preferences}

\IZ
\itemi The person's utility in period 0 from consuming $C_0$ and $C_1$ is given by:

\begin{align}
U(C_0,C_1) & = \log C_0 + \beta \log C_1,
\end{align}

\

\itemii $U$ is an \emph{intertemporal utility function} and $\log C$ is the \emph{flow} of utility generated each period.

\

\itemiii $0\leqslant\beta\leqslant 1$ is the person's \emph{subjective discount factor}.

\

%The utility function implies indifference curves of the form:
%
%\begin{align}
%C_1 & =  e^{ \frac{\bar{U} - \log C_1}{\beta}},
%\end{align}
\ZI
}

\fr{
\ft{Preferences}

\IZ
\itemi The utility function implies indifference curves of the form:

\begin{align}
C_1 & =  e^{ \frac{\bar{U} - \log C_0}{\beta}},
\end{align}

where $\bar{U}$ is some given level of utility.

\

\itemi As long as $0<\beta<1$, the indifference curves are convex to the origin.
\ZI
}




\fr{
\ft{Budget Constraints and the Boundary Condition}

\IZ
\itemi Let $K_t$ denote the quantity of cake available in period $t$  with $K_0$ given. 

\

\itemii The cake neither grows nor depreciates over time so the person faces the following budget constraints each period:

\begin{align}
C_0 & =  K_0 - K_1\\
C_1 & =  K_1 - K_2.
\end{align}
\ZI

}


\fr{
\ft{Budget Constraints and the Boundary Condition}

\IZ
\item A \emph{boundary condition} imposes a limit on the endpoint of the problem. 

\

\itemii Since the person only consumes cake in periods 0 and 1, there is no value in leaving leftover cake at the end of period 1 which implies the boundary condition:

\begin{align}
K_2 & = 0.
\end{align}

\ZI
}



\fr{
\ft{Budget Constraints and the Boundary Condition}


\IZ
\itemi The boundary condition allows us to eliminate $K_2$ from the period 1 budget constraint:

\begin{align}
C_0 & =  K_0 - K_1\\
C_1 & =  K_1.
\end{align}


\itemii Combine the two budget constraints to eliminate $K_1$ and obtain an \emph{intertemporal budget constraint}:

\begin{align}
C_1 & =  K_0 - C_0.
\end{align}

\

\itemiii Notice that the implied price of $C_1$ in terms of $C_0$ is 1.

\ZI
}





\fr{
\ft{Optimization}

\IZ
\itemi The person's problem is to choose $C_0$, $C_1$, $K_1$, and $K_2$ to maximize his utility subject to the two budget constraints and the boundary condition. 

\

\itemii However, the problem can be simplified:

	\EN
	\itemiii Use the boundary condition to eliminate $K_2$ from the budget constraints. 
	\itemiv Then use the budget constraints to substitute $C_0$ and $C_1$ out of the utility function. 
	\NE

\ZI
}


\fr{
\ft{Optimization}

\IZ
\itemi The result is an unconstrained optimization problem in $K_1$.

\begin{align}
\max_{K_1} \; \log (K_0 - K_1) + \beta \log K_1,
\end{align}

\

\itemii Take the derivative with respect to $K_1$ and set the derivative equal to zero to obtain the \emph{first-order} or \emph{optimality} condition:

\begin{align}
-\frac{1}{K_0 - K_1} + \frac{\beta}{K_1} & = 0.
\end{align}
\ZI
}






\fr{
\ft{Optimization}

\IZ
\itemi Solve for $K_1$ to find the optimal cake holding for period 1:

\begin{align}
\boxed{K_1  = \frac{\beta}{1-\beta}K_0}
\end{align}

\

\itemii Cake in period 1 is proportional to initial cake value.
\ZI

}


\fr{
\ft{Optimization}

\IZ
\itemi Then use the budget constraints to infer optimal cake consumption each period:
\begin{align}
\boxed{C_0  = \frac{1}{1+\beta}K_0}
\end{align}

and:

\begin{align}
\boxed{C_1  = \frac{\beta}{1+\beta}K_0}.
\end{align}

\

\itemii Notice that $C_0 + C_1 = K_0$ so all of the cake is eaten by the end.


\ZI
}
















\end{document}